% ====================================================================
% Faithful Self-Modeling Splits h_3(O) into Observable and Beable
% Paper 1 (culmination of Papers 5, 7, 6).
% ====================================================================

\documentclass[12pt]{article}

% --- Packages ---
\usepackage{amsmath,amssymb,amsthm}
\usepackage[margin=1in]{geometry}
\usepackage{graphicx}
\usepackage{booktabs}
\usepackage{hyperref}
\usepackage{xcolor}
\usepackage{enumitem}
\usepackage{microtype}
\usepackage{float}

% --- Hyperref ---
\hypersetup{
  colorlinks=true,
  linkcolor=blue!60!black,
  citecolor=green!40!black,
  urlcolor=blue!70!black
}

% --- Theorem environments ---
\newtheorem{theorem}{Theorem}
\newtheorem{proposition}[theorem]{Proposition}
\newtheorem{lemma}[theorem]{Lemma}
\newtheorem{corollary}[theorem]{Corollary}
\newtheorem{definition}[theorem]{Definition}
\newtheorem{postulate}[theorem]{Postulate}
\newtheorem{conjecture}[theorem]{Conjecture}
\theoremstyle{remark}
\newtheorem{remark}[theorem]{Remark}

% --- Notation shortcuts ---
\newcommand{\hthree}{h_3(\mathbb{O})}
\newcommand{\Ffour}{F_4}
\newcommand{\rhoJ}{\rho_J}
\newcommand{\Phiexp}{\Phi}
\newcommand{\Stream}{\mathcal{X}}
\newcommand{\Aut}{\mathrm{Aut}}
\newcommand{\Tr}{\mathrm{Tr}}
\newcommand{\Jnorm}[1]{\lVert #1 \rVert_J}

% --- Title ---
\title{Faithful Self-Modeling Splits $\hthree$ into Observable and Beable}
\author{Bryan Ehrlich}
\date{}

\begin{document}
\maketitle

% ====================================================================
% ABSTRACT
% ====================================================================

\begin{abstract}
The exceptional Jordan algebra $\hthree$, with Jordan product
$X \circ Y = \tfrac{1}{2}(XY + YX)$, automorphism group
$\Ffour = \Aut\hthree$, Jordan norm $\Jnorm{Y}^2 = \Tr(Y^2)$, and
state space $\Omega = \{X \in \hthree : X \succeq 0,\ \Tr X = 1\}$,
carries a normalized self-modeling iteration $\phi(X) = X^2/\Tr(X^2)$.
Four standard mathematical objects attach to a self-modeler on
$\hthree$: a \textbf{Kind} ($\Ffour$-orbit $[X] \in \Omega/\Ffour$), a
\textbf{State} ($X \in \Omega$), a \textbf{Stream} (a $\phi$-iterate
$(X_0, X_1, X_2, \ldots)$ with $X_{k+1} = \phi(X_k)$), and an
\textbf{Observable} (an $\Ffour$-invariant polynomial on $\hthree$,
equivalently a polynomial in the three Faraut-Kor\'anyi
generators~\cite{faraut-koranyi1994} $\Tr X$, $\Tr(X^2)$, $\det X$).
Under Bell's observable/beable binary~\cite{bell1984speakable}, Kind
and Observable lie on the Observable side; State and Stream on the
Beable side.

Observable projection $X \mapsto [X]$ and Stream propagation $\phi$
do not commute: Observable extraction destroys the $\Ffour$-gauge
degrees of freedom the next $\phi$-step would have consumed. In the
autonomous case ($\phi$ $\Ffour$-equivariant) the non-commutation
reduces to a residual global gauge
(Theorem~\ref{thm:autonomous-comp}). In the non-autonomous case the
pairwise diagonal invariant $\Tr(X_k \cdot X_{k+1})$ belongs to a
strict super-ring of the pointwise $\Ffour$-invariants, yielding
strong-form complementarity (Proposition~\ref{prop:strong-comp}).
The resulting Observable-Beable exhaustion
(Corollary~\ref{cor:OB1}; the two information modes jointly span
the $\Ffour$-invariant content of the Stream) realizes inside
$\hthree$ by algebra the observable/beable split Bell described for
quantum mechanics and Barandes~\cite{barandes2023stochquantum}
reformulated stochastically.

The $\Ffour$-invariant measure on Streams that carries the residual
Beable content is
\[
  \Phiexp[\Stream] \;=\; \sum_k \rhoJ(X_k)\,\Jnorm{X_{k+1}-X_k}^2,
\]
where $\rhoJ(X) = \det(X)(\Tr(X^2) - \tfrac{1}{3})$ is the product
of the lowest-degree $\Ffour$-invariants vanishing on the
rank-deficient stratum $\partial\Omega$ and at the maximally mixed
state $[I/3]$. A finite-dimensional enumeration on pair-invariants
(Theorem~\ref{thm:phi-family}) followed by a path-space purity
projection (Corollary~\ref{cor:phi-canonical}) selects $\Phiexp$
at minimal polynomial degree on pairs, up to a positive scalar.
The uniqueness inherits the factored-form hypothesis of
Theorem~\ref{thm:rho-uniqueness} for $\rhoJ$; an unconditional
(factorization-free, unbounded-degree) $\rhoJ$ uniqueness is
flagged as open in Section~\ref{sec:gaps} (G4).
\end{abstract}

% ====================================================================
% SECTION 1: INTRODUCTION
% ====================================================================

\section{Introduction}
\label{sec:intro}

Paper~5~\cite{ehrlich2026qm} derives the Euclidean Jordan structure on
$\hthree$ from faithful self-modeling on a spectral order-unit space,
and Paper~7~\cite{ehrlich2026h3o} identifies $\hthree$ as the
non-composable case; together they leave $\hthree$ equipped with the
normalized iteration $\phi(X) = X^2/\Tr(X^2)$, its automorphism group
$\Ffour$, and the state space $\Omega$, but with no measure on
$\phi$-trajectories and no named ontology for the trajectories.
Faithful self-modeling, in Paper~5's precise sense, is a
four-clause condition on a finite-dimensional spectral order-unit
space: a faithful order-isomorphism body-to-model tracking map,
minimal product composite, and simplicity
(Appendix~\ref{app:self-modeling}, reproducing Paper~5's
Definition~1). In this paper it is used only as the premise
fixing $(\hthree, \Ffour, \Omega, \phi)$ as a package, with
$\phi$'s specific form $X \mapsto X^2/\Tr(X^2)$ taken as Paper~5's
output rather than re-derived here. This paper supplies the
measure and the ontology.

\begin{theorem}[Main, informal]\label{thm:main-informal}
On $\hthree$-Streams, four standard mathematical objects
(Kind, State, Stream, Observable; Section~\ref{sec:ontology}) split
under Bell's observable/beable binary into two Observable-side
objects (Kind, Observable) and two Beable-side objects (State,
Stream). Observable projection and Stream propagation do not
commute: in the autonomous case this reduces to a residual global
gauge (Theorem~\ref{thm:autonomous-comp}); in the non-autonomous
case, path-space carries content strictly beyond any pointwise
Observable readout
(Proposition~\ref{prop:strong-comp};
Corollary~\ref{cor:OB1}, Observable-Beable exhaustion). The
$\Ffour$-invariant functional on Streams that carries the residual
Beable content is
\begin{equation}\label{eq:phi-intro}
  \Phiexp[\Stream] \;=\; \sum_k \rhoJ(X_k)\,\Jnorm{X_{k+1} - X_k}^2,
\end{equation}
with $\rhoJ(X) = \det(X)\,(\Tr(X^2) - \tfrac{1}{3})$. Among
minimal-degree diagonal-$\Ffour$-invariant pair functionals
satisfying five formal conditions (P1-P5) and a non-negativity
constraint, $\Phiexp$ is singled out up to a positive scalar by a
path-space purity projection
(Theorem~\ref{thm:phi-family}, Corollary~\ref{cor:phi-canonical});
the uniqueness inherits the factored-form hypothesis of
Theorem~\ref{thm:rho-uniqueness} for $\rhoJ$.
\end{theorem}

\noindent
The formal version reduces to a finite-dimensional enumeration of
minimal-degree $\Ffour$-invariants on pairs of states, decomposed
into non-negativity blocks. The paper develops the ontology
(Section~\ref{sec:ontology}), constructs $\rhoJ$ as the
Observable-side weight (Section~\ref{sec:rho-obs}), constructs
$\Phiexp$ as the Beable-side functional and proves its uniqueness
(Section~\ref{sec:phi-def}), states the complementarity results
(Section~\ref{sec:complementarity}), records temporal additivity
(Section~\ref{sec:consequences}), lists empirical signatures
(Section~\ref{sec:signatures}), and closes with related work and
open gaps (Sections~\ref{sec:related}-\ref{sec:gaps}).
Appendix~\ref{app:self-modeling} reproduces the standalone
self-modeling definition of Paper~5 for readers who prefer not to
chase the pointer.

\paragraph{Scope.}
This paper is not a derivation of $\hthree$ (Paper~7) or of the
self-modeling map $\phi$ (Paper~5); both are taken as given. Results
are at $\Ffour$-orbit level, and no biological or computational
encoding into $\hthree$ is committed to (open gap G1). Temporal
additivity of $\Phiexp$ is proved; spatial aggregation is a
conjecture flagged wherever used. The word ``measure'' is used in the
physicist's sense of a weight on trajectories, not a $\sigma$-additive
set function in the strict measure-theoretic sense.

% ====================================================================
% SECTION 2: THE FOUR-TYPE ONTOLOGY
% ====================================================================

\section{Kind, State, Stream, Observable}
\label{sec:ontology}

The ontology the rest of the paper uses is four standard mathematical
objects attached to an $\hthree$-self-modeler, organized under Bell's
observable/beable binary.

\begin{definition}[The four types]\label{def:ontology}
For a self-modeler in $\hthree$ with state $X \in \Omega$ evolving
under the self-modeling map $\phi(X) = X^2/\Tr(X^2)$ of Paper~5:
\begin{enumerate}[label=(\roman*)]
  \item\label{type:kind} A \textbf{Kind} is an $\Ffour$-orbit
    $[X] = \{g \cdot X : g \in \Ffour\} \in \Omega/\Ffour$.
  \item\label{type:state} A \textbf{State} is a point $X \in \Omega$.
  \item\label{type:stream} A \textbf{Stream} is a $\phi$-iteration
    trajectory $\Stream = (X_0, X_1, X_2, \ldots)$ with
    $X_{k+1} = \phi(X_k)$.
  \item\label{type:observable} An \textbf{Observable} is an
    $\Ffour$-invariant polynomial
    $R \in \mathbb{R}[\hthree]^{\Ffour}$, equivalently a polynomial
    in $(\Tr X, \Tr(X^2), \det X)$ by
    Faraut-Kor\'anyi~\cite{faraut-koranyi1994}.
\end{enumerate}
The four types organize as a $2\times 2$ under Bell's
\emph{observable}/\emph{beable}
binary~\cite{bell1984speakable}: Kind and Observable lie on the
observable side (quantities a third party can read; Kind is the
Observable image of a State); State and Stream lie on the beable
side (elements of physical reality intrinsic to the system). Our
``beable'' is Bell's but with State and $\phi$-iterates in place of
classical field or particle values; the trajectory-as-beable shape
is shared with Barandes's~\cite{barandes2023stochquantum}
indivisible-stochastic reformulation of quantum mechanics, from which
we borrow the shape, not the stochastic machinery. ``Observable''
is used in its gauge-theoretic sense throughout (a function constant
along $\Ffour$-orbits), not as a synonym for a self-adjoint operator;
``Kind'' avoids collision with type-theoretic vocabulary.
\end{definition}

\subsection{Kind: the orbit}
\label{sec:kind}

The \emph{Kind} of $X \in \Omega$ is its $\Ffour$-orbit
\begin{equation}\label{eq:kind}
  [X] \;=\; \{\,g \cdot X : g \in \Ffour\,\} \;\in\; \Omega / \Ffour.
\end{equation}
Every $\Ffour$-invariant property of $X$ factors through $[X]$; by
Faraut-Kor\'anyi~\cite{faraut-koranyi1994} these properties are
polynomials in $(\Tr X, \Tr(X^2), \det X)$. The Kind is the
Observable-side universal: ``red'' as a type is a Kind.

\paragraph{Prior work.} $\Ffour$-orbits on $\hthree$ have been
catalogued since Jacobson 1968 and appear in Springer-Veldkamp
\cite{springer-veldkamp2000}. The orbit space $\Omega / \Ffour$ is a
semi-algebraic subset of $\mathbb{R}^3$ coordinatized by the three
invariants.

\subsection{State: the instance}
\label{sec:state}

A \emph{State} is a point $X \in \Omega$. It is a Kind plus a choice of
$\Ffour$-frame: one Kind is realized by the continuum of States in its
orbit, and the selection of representative is a gauge choice, not a
physical one. $\Ffour$-invariant measurements see only the Kind, but the
self-modeling map $\phi$ acts on States
(\S\ref{sec:stream}).

\paragraph{Prior work.} Normalized positive elements as states of a
quantum system go back to von~Neumann 1932. For Jordan algebras, the
state space is the base of the positive cone; the Koecher-Vinberg
theorem (1957-1962) gives the bijection between Euclidean Jordan
algebras and homogeneous self-dual cones. Treatments:
Alfsen-Shultz~\cite{alfsen-shultz2001},
Bengtsson-\.Zyczkowski~\cite{bengtsson-zyczkowski2006}; exceptional
case in Barnum-Graydon-Wilce~\cite{barnum-graydon-wilce2015}.

\subsection{Stream: the dynamics}
\label{sec:stream}

Paper~5's self-modeling is a one-step map $\phi$ on $\Omega$,
\begin{equation}\label{eq:phi-iteration}
  X_{k+1} \;=\; \phi(X_k), \qquad k = 0, 1, 2, \ldots
\end{equation}
$\Ffour$-equivariant in the autonomous case, drive-indexed
$X_{k+1} = \phi_k(X_k)$ in the non-autonomous case. Iterating from
$X_0 \in \Omega$ produces the \emph{Stream}
\begin{equation}\label{eq:stream}
  \Stream \;=\; (X_0, X_1, X_2, \ldots).
\end{equation}
The iteration index $k$ is \emph{$\phi$-time}, algebraic and internal;
its encoding into wall-clock time for a particular physical
realization is a framework task, not addressed here.

\paragraph{Prior work.} Iteration sequences are the primordial object
of dynamical systems; standard treatments
in~\cite{devaney1989, katok-hasselblatt1995}. Hofstadter's ``strange
loop''~\cite{hofstadter1979, hofstadter2007} names the same object
from a cognitive direction; the present framework supplies the
specific substrate $\hthree$ and the specific map $\phi$ from Paper~5.

\subsection{Observable: the invariant}
\label{sec:observable}

An \emph{Observable} is an $\Ffour$-invariant polynomial on $\hthree$,
\begin{equation}\label{eq:observable-space}
  R \in \mathbb{R}[\hthree]^{\Ffour}.
\end{equation}
By Faraut-Kor\'anyi~\cite{faraut-koranyi1994} the ring of such
polynomials is freely generated by the three algebraically independent
scalars $\Tr X$, $\Tr(X^2)$, $\det X$. Under the normalization
$\Tr X = 1$, every Observable on $\Omega$ factors through the
two-dimensional plane coordinatized by $(\Tr(X^2), \det X)$. The name
matches the gauge-theoretic usage: $\Ffour$ is the gauge group of
$\hthree$ and an Observable is a gauge-invariant scalar. Any external
measurement of $X$ returns values of such polynomials - at most three
real numbers - while a State carries continuously many
$\Ffour$-gauge degrees of freedom beyond the Kind and a Stream an
open-ended sequence of them.

\paragraph{Prior work.} Classical invariant theory (Cayley, Sylvester,
Hilbert); Weyl~\cite{weyl1939}. The free-generator theorem for $\Ffour$
on $\hthree$ is Faraut-Kor\'anyi Ch.~III.

\subsection{A worked example}
\label{sec:worked}

To make the four types concrete, fix a state $X \in \Omega$ with
eigenvalues $(\lambda_1, \lambda_2, \lambda_3)$, $\sum \lambda_i = 1$,
$\lambda_i \geq 0$. Then:
\begin{itemize}[leftmargin=2em]
\item The \textbf{Kind} of $X$ is its $\Ffour$-orbit, fully determined by
the unordered triple $(\lambda_1, \lambda_2, \lambda_3)$, equivalently by
$(\Tr X, \Tr X^2, \det X) = (1,\,\sum \lambda_i^2,\,\prod \lambda_i)$.
\item The \textbf{State} is $X$ itself: a specific element of the orbit.
Two states with the same eigenvalues but in different $\Ffour$-frames are
different States of the same Kind.
\item The \textbf{Stream} starting from $X$ under $\phi$ is the sequence
$(X, \phi(X), \phi^2(X), \ldots)$. For an autonomous $\Ffour$-equivariant
$\phi$, the Stream's Kind sequence is fully determined by the three
invariants at each step; the frame-level trajectory is not.
\item An \textbf{Observable} reading of $X$ is a polynomial in $(\Tr X,
\Tr X^2, \det X)$; for example, the self-reference density $\rhoJ$ of
Section~\ref{sec:rho-obs} is one such polynomial.
\end{itemize}

% ====================================================================
% SECTION 3: rho ON THE OBSERVABLE SIDE
% ====================================================================

\section{The self-reference density on the Observable side}
\label{sec:rho-obs}

The Observable-side weight on self-modeling is
\begin{equation*}
  \rhoJ(X) \;=\; \det(X)\,(\Tr(X^2) - \tfrac{1}{3}),
\end{equation*}
a polynomial in the three Faraut-Kor\'anyi
generators~\cite{faraut-koranyi1994} and therefore an Observable by
Definition~\ref{def:ontology}; it supplies the instantaneous weight
that the Beable-side $\Phiexp$ of Section~\ref{sec:phi-def} will
multiply into the running term. The specific form emerges from
$\phi$'s fixed-point structure: the three $\phi$-fixed-point orbits
$[I/3]$, $[(\tfrac{1}{2}, \tfrac{1}{2}, 0)]$, $[(1,0,0)]$ are the
states on which self-modeling is not live, so an Observable-side
weight must vanish on all three. The rank-deficient orbits
$[(\tfrac{1}{2}, \tfrac{1}{2}, 0)]$ and $[(1,0,0)]$ sit in
$\{\det = 0\}$, and $[I/3]$ sits at $\{\Tr(X^2) = \tfrac{1}{3}\}$;
$\rhoJ$ is the product of the lowest-degree $\Ffour$-invariants
cutting out these two conditions.
Lemmas~\ref{lem:det-forced} and~\ref{lem:t2-forced} establish
factor-level uniqueness; Theorem~\ref{thm:rho-uniqueness} combines
them under the factorization hypothesis;
Remark~\ref{rem:rho-alternatives} records the factorization-free Lean
result and the remaining unconditional gap (G4).

\paragraph{Degree conventions.} Throughout Section~\ref{sec:rho-obs},
\emph{polynomial degree} means degree in the matrix entries of
$\hthree$: $\Tr$ has degree $1$, $\Tr(X^2)$ degree $2$, $\det$ degree
$3$. The Faraut-Kor\'anyi \emph{generator degree} counts each of
$(\Tr, \Tr(X^2), \det)$ as a single generator; $\rhoJ = \det \cdot
(\Tr(X^2) - \tfrac{1}{3})$ has generator degree $2$ and matrix-entry
degree $5$. The Lean result (Remark~\ref{rem:rho-alternatives}) uses
a third notion, \emph{eigenvalue degree}, which is matrix-entry
degree restricted to symmetric functions of the three eigenvalues
$(\lambda_1, \lambda_2, \lambda_3)$. All minimum-degree claims below
are in matrix-entry degree unless noted.

\subsection{The three invariants on the state slice}
\label{sec:FK-slice}

On $\Omega$ the trace invariant is fixed at $\Tr X = 1$, and the full
$\Ffour$-invariant ring on the slice is generated by the two remaining
scalars
\begin{equation}\label{eq:slice-invariants}
  t_2 \;=\; \Tr(X^2), \qquad t_3 \;=\; \det X,
\end{equation}
with $t_2 \in [\tfrac{1}{3}, 1]$ and $t_3 \in [0, \tfrac{1}{27}]$ on
$\Omega$, by Cauchy-Schwarz applied to the three eigenvalues of $X$. The
three fixed-point orbits are pinned in these coordinates:
\begin{equation}\label{eq:fixed-points}
  [I/3]: \ t_2 = \tfrac{1}{3},\ t_3 = \tfrac{1}{27}; \qquad
  \bigl[(\tfrac{1}{2}, \tfrac{1}{2}, 0)\bigr]: \ t_2 = \tfrac{1}{2},\ t_3 = 0; \qquad
  [(1, 0, 0)]: \ t_2 = 1,\ t_3 = 0.
\end{equation}

What is the lowest-degree $\Ffour$-invariant polynomial on
$\hthree$ that vanishes on all three orbits and is not identically
zero elsewhere on $\Omega$?

\subsection{Uniqueness at degree three and degree two}
\label{sec:rho-uniqueness}

Two sub-uniqueness results combine.

\begin{lemma}[Rank-deficiency detector]\label{lem:det-forced}
The unique (up to real scalar) $\Ffour$-invariant polynomial on $\hthree$
of minimum degree vanishing on $\{X \in \hthree : \det X = 0\}$ is
$f = c \cdot \det$, at degree three.
\end{lemma}

\begin{proof}[Proof sketch]
Every $\Ffour$-invariant polynomial on $\hthree$ is a polynomial in
$(t_1, t_2, t_3)$ \cite{faraut-koranyi1994}. The ideal of polynomials
vanishing on $\{t_3 = 0\}$ is the principal ideal $(t_3)$, because $t_3$
is irreducible in $\mathbb{R}[t_1, t_2, t_3]$ and generates a prime ideal.
Any polynomial in this ideal is of the form $t_3 \cdot h$ for some
$h$; minimum degree is achieved at $h = $ constant, giving $f = c t_3$.
\end{proof}

\begin{lemma}[$I/3$-deviation detector]\label{lem:t2-forced}
Restricted to the slice $\{t_1 = 1\}$, the unique (up to real scalar)
$\Ffour$-invariant polynomial of minimum degree that vanishes at $X = I/3$
and is not identically zero on the rank-one orbit $[(1, 0, 0)]$ is
$g = c \cdot (t_2 - \tfrac{1}{3})$, at degree two.
\end{lemma}

\begin{proof}[Proof sketch]
The restriction of $\mathbb{R}[t_1, t_2, t_3]$ to $\{t_1 = 1\}$ is
$\mathbb{R}[t_2, t_3]$. At degrees zero and one, the only restrictions are
constants, which cannot vanish at $I/3$ without being identically zero. At
degree two, the ambient basis is $(t_1^2, t_2)$, restricting to $(1, t_2)$;
a general element is $a t_2 + b$. Vanishing at $X = I/3$ (where
$t_2 = \tfrac{1}{3}$) forces $b = -a/3$, giving $g = a(t_2 - \tfrac{1}{3})$.
Non-vanishing on $[(1,0,0)]$ (where $t_2 = 1$) forces $a \neq 0$. Hence
$g = c (t_2 - \tfrac{1}{3})$ is unique up to scalar at minimum degree.
\end{proof}

\begin{theorem}[Structural selection of $\rhoJ$]\label{thm:rho-uniqueness}
Up to a positive scalar, the unique $\Ffour$-invariant polynomial on
$\hthree$ of the form $p \cdot q$, with $p$ as in
Lemma~\ref{lem:det-forced} and $q$ as in Lemma~\ref{lem:t2-forced}, is
\begin{equation}\label{eq:rho-J-definition}
  \rhoJ(X) \;=\; \det(X)\,\bigl(\Tr(X^2) - \tfrac{1}{3}\bigr).
\end{equation}
$\rhoJ$ is a product of $\Ffour$-invariants of matrix-entry degrees
$3$ and $2$, and has degree $2$ in the generator ring
$\mathbb{R}[\Tr, \Tr(X^2), \det]$. Its vanishing locus on $\Omega$ is
the rank-deficient stratum $\partial\Omega$ together with the
maximally mixed state $[I/3]$; its positivity locus is the open
interior $\Omega^{\circ} \setminus [I/3]$, on which self-modeling is
live.
\end{theorem}

\begin{proof}[Proof sketch]
Combining Lemmas~\ref{lem:det-forced} and~\ref{lem:t2-forced},
$p \cdot q = c_1 c_2 \cdot \det \cdot (t_2 - \tfrac{1}{3})$, which is
$c \cdot \rhoJ$ for $c > 0$. The vanishing/positivity claims follow from
eigenvalue positivity: on $\Omega$ the eigenvalues satisfy
$\lambda_i \geq 0$ with $\sum \lambda_i = 1$, so $\det X = \prod \lambda_i$
vanishes exactly when at least one eigenvalue is zero (the rank-deficient
orbits), and $t_2 - \tfrac{1}{3} = \sum \lambda_i^2 - \tfrac{1}{3}$ vanishes
exactly when all three are $\tfrac{1}{3}$ by Cauchy-Schwarz, i.e.\ at
$[I/3]$. On the complement, both factors are strictly positive.
\end{proof}

\begin{remark}[Relation to Lean-verified uniqueness]\label{rem:rho-alternatives}
Theorem~\ref{thm:rho-uniqueness} is unbounded in polynomial degree
but assumes the factored shape $p \cdot q$ with $p, q$ the two
minimal-degree factors of Lemmas~\ref{lem:det-forced}
and~\ref{lem:t2-forced}. A related uniqueness statement has been
machine-verified: the Lean theorem
\texttt{rhoJ\_unique\_minimal}~\cite{lean-rhoj} shows that every
$\Ffour$-invariant polynomial on $\hthree$ of eigenvalue degree
$\leq 5$ that vanishes on $\partial\Omega$ and at $[I/3]$ equals
$c \cdot \rhoJ$, without assuming factorization. Neither statement
strictly dominates the other: the in-paper result drops the
bounded-degree hypothesis but assumes factorization; the Lean result
drops factorization but caps eigenvalue degree at $5$. The
unconditional statement - drop factorization, keep unbounded degree -
is open (Section~\ref{sec:gaps}, G4).
\end{remark}

% ====================================================================
% SECTION 4: THE MEASURE PHI ON BEABLE-TRAJECTORIES
% ====================================================================

\section{The measure \texorpdfstring{$\Phiexp$}{Phi} on Beable-trajectories}
\label{sec:phi-def}

Section~\ref{sec:rho-obs} selected $\rhoJ$ as the Observable-side
weight - a single-state scalar, computable from any outside
party's pointwise $\Ffour$-invariant readout. The Beable-side
object is the Stream itself, and the functional that captures the
pair-level content an Observable-side readout cannot recover is
$\Phiexp$.

\subsection{Definition}
\label{sec:phi-definition}

\begin{definition}[the $\Phiexp$-measure]\label{def:phi}
Given a Stream $\Stream = (X_0, X_1, \ldots, X_K)$ of a self-modeler in
$\hthree$, define the $\Phiexp$-measure on $\Stream$ by
\begin{equation}\label{eq:phi-definition}
  \Phiexp[\Stream] \;=\; \sum_{k=0}^{K-1} \rhoJ(X_k)\,\Jnorm{X_{k+1} - X_k}^2,
\end{equation}
where $\Jnorm{Y}^2 = \Tr(Y^2)$ is the Jordan squared norm. The
continuous-time (phi-time) limit is
\begin{equation}\label{eq:phi-continuous}
  \Phiexp[\Stream] \;=\; \int_0^{T_\phi} \rhoJ(X(t_\phi))\,
    \Big\lVert \frac{\mathrm{d}X}{\mathrm{d}t_\phi} \Big\rVert_J^2 \,\mathrm{d}t_\phi.
\end{equation}
\end{definition}

The integration is over \emph{phi-time} - the iteration index $k$,
which is the self-modeler's own clock - not wall-clock time. An outside
observer running an atomic clock while $S$ self-models sees wall-clock
time; $S$ lives in phi-time. The two are not in general commensurate.
The summand factors into a \emph{weight} $\rhoJ(X_k)$ (how densely
self-referential the current State is) and a \emph{running} term
$\Jnorm{X_{k+1} - X_k}^2$ (how vigorously the Stream is moving at
step $k$).

Expanding the running term in Jordan form,
\begin{equation}\label{eq:running-expansion}
  \Jnorm{X_{k+1} - X_k}^2 \;=\; \Tr(X_k^2) + \Tr(X_{k+1}^2) - 2\,\Tr(X_k \cdot X_{k+1}),
\end{equation}
reveals the pairwise invariant $\Tr(X_k \cdot X_{k+1})$, which is
$\Ffour$-invariant under the diagonal action
$(X_k, X_{k+1}) \mapsto (g X_k, g X_{k+1})$ but is not in general
computable from the pointwise Kind sequence $([X_0], [X_1], \ldots)$.
This pair-level invariant is the object that carries $\Phiexp$'s
path-space content; we return to it in Section~\ref{sec:complementarity}.

\subsection{Five formal properties}
\label{sec:phi-properties}

\begin{itemize}[leftmargin=2em]
\item \textbf{(P1) Dead Streams vanish.} If $X_{k+1} = X_k$ for every
$k$ (a frozen Stream), then $\Phiexp[\Stream] = 0$ term-by-term from
$\Jnorm{X_{k+1} - X_k} = 0$.

\item \textbf{(P2) Zero-weight Streams vanish.} If $\rhoJ(X_k) = 0$
for every $k$ - Stream confined to the $\rhoJ$-vanishing locus of
Section~\ref{sec:rho-uniqueness} - then $\Phiexp[\Stream] = 0$
term-by-term. Example: an oscillation between two rank-one projectors
runs, but carries no self-reference structure.

\item \textbf{(P3) Diagonal $\Ffour$-invariance.} For any $g \in \Ffour$
and any Stream $\Stream$,
$\Phiexp[g \cdot \Stream] = \Phiexp[\Stream]$, where
$g \cdot \Stream = (gX_0, gX_1, \ldots)$. The proof combines
$\Ffour$-invariance of $\rhoJ$ with the isometry property of the Jordan
norm under $\Ffour$-action.

\item \textbf{(P4) Non-reducibility to pointwise Observables.} There exist
pairs of Streams $\Stream_a, \Stream_b$ with identical pointwise Kind
sequences ($[X_k^{(a)}] = [X_k^{(b)}]$ for all $k$) but
$\Phiexp[\Stream_a] \neq \Phiexp[\Stream_b]$. $\Phiexp$ is not a
function of the pointwise Kind sequence alone.

\item \textbf{(P5) Internal to the algebra.} Every quantity in
Definition~\ref{def:phi} is a polynomial in $\hthree$-data: $\rhoJ$ has
degree $5$ in matrix entries (bilinear in the generators $\det$ and
$\Tr(X^2)$), $\Jnorm{X_{k+1} - X_k}^2$ is degree $2$, so each per-tick
contribution is a degree-$7$ polynomial in $\hthree$ entries. No
wall-clock time, coordinate chart, or outside observer is invoked. The
measure lives in $\hthree$ alone.
\end{itemize}

All five properties have been verified numerically in the companion
script \cite{phi-numerics}. (P4) has a sharp numerical realization.
Consider two Streams constructed with identical
pointwise Kind sequences - one autonomous (gauge fixed by the dynamics),
one non-autonomous with a cyclic gauge schedule
$\sigma^{k \bmod 3} \in \Ffour$. Pointwise Kinds match to
$10^{-16}$ (machine precision), while $\Phiexp$ differs by a factor of
twenty-one: $\Phiexp_a = 1.51 \times 10^{-4}$ and
$\Phiexp_b = 3.27 \times 10^{-3}$. The ratio lives in the pairwise
invariant $\Tr(X_k \cdot X_{k+1})$; it is invisible to any outside
party who measures only pointwise Kinds.

\subsection{Uniqueness}
\label{sec:uniqueness}

The five properties above do not yet determine $\Phiexp$. One can ask:
what is the complete family of functionals on pairs of $\hthree$-elements
satisfying P1-P5 and a non-negativity condition, and what additional
condition singles out Definition~\ref{def:phi} uniquely?

\begin{theorem}[Minimal-degree family]\label{thm:phi-family}
Let $f : \hthree \times \hthree \to \mathbb{R}$ be an $\Ffour$-invariant
polynomial (under the diagonal $\Ffour$-action) of minimal total
polynomial degree satisfying
\begin{itemize}[leftmargin=2em]
\item $f(X, X) = 0$ for all $X$, \hfill (reduction of P1)
\item $f(X, Y) = 0$ whenever $\rhoJ(X) = 0$, \hfill (reduction of P2)
\item $f$ depends non-trivially on $\Tr(X \cdot Y)$, \hfill (reduction of P4)
\item $f \geq 0$ on the physical state cone,\hfill (non-negativity)
\end{itemize}
and polynomial in algebra components (P5). Then $f$ is a non-negative
combination
\begin{equation}\label{eq:phi-family}
  f(X, Y) \;=\; \rhoJ(X)\Bigl[\,\beta\,(\Tr X - \Tr Y)^2 \;+\; \delta\,\Jnorm{X - Y}^2\,\Bigr]
\end{equation}
with $\beta \geq 0$ and $\delta > 0$. The minimal total polynomial degree
is seven, and the family is two-parameter.
\end{theorem}

\begin{proof}
The proof is in six steps. All invariants below are diagonal
$\Ffour$-invariants on $\hthree \times \hthree$; ``degree''
means total matrix-entry degree (Section~\ref{sec:rho-obs} degree
conventions).

\textbf{Step 1 (P2 factorization in the pair ring).} The claim is
that (P2) forces $f(X, Y) = \rhoJ(X) \cdot g(X, Y)$ for some
polynomial $g \in \mathbb{R}[\hthree \times \hthree]$. The
argument has two parts: (1a) $(\rhoJ)$ is a radical ideal in the
single-copy polynomial ring $\mathbb{R}[\hthree]$; (1b) vanishing
of $f(X, Y)$ on $V(\rhoJ) \times \hthree$ (as a subset of
$\hthree \times \hthree$) lifts to divisibility by $\rhoJ(X)$ in
the pair ring.

\emph{(1a) Radicality in the single-copy ring.} The polynomial
$\det \in \mathbb{R}[\hthree]$ is irreducible (the standard cubic
$\Ffour$-invariant; see Faraut-Kor\'anyi Ch.~V), and
$\Tr(X^2) - \tfrac{1}{3}$ is irreducible (a degree-$2$ polynomial
with nonvanishing gradient). They are coprime: their zero loci
$V(\det)$ and $V(\Tr(X^2) - \tfrac{1}{3})$ are distinct irreducible
hypersurfaces in $\hthree$ (different degrees; $\det$ vanishes on
rank-deficient states, $\Tr(X^2) - \tfrac{1}{3}$ on the
trace-$1$-shifted second-invariant level set). Hence $\rhoJ =
\det \cdot (\Tr(X^2) - \tfrac{1}{3})$ is a product of two distinct
irreducible factors, so $(\rhoJ)$ is a radical ideal in
$\mathbb{R}[\hthree]$, and Hilbert's Nullstellensatz gives
$I(V(\rhoJ)) = (\rhoJ)$.

\emph{(1b) Lifting vanishing to divisibility in the pair ring.}
Expand $f(X, Y)$ in monomials in the matrix entries of $Y$:
\[
  f(X, Y) \;=\; \sum_{\alpha} c_\alpha(X) \, m_\alpha(Y),
\]
where $\{m_\alpha(Y)\}$ is a linearly independent monomial basis
and $c_\alpha \in \mathbb{R}[\hthree]$. Condition (P2) says
$f(X, Y) = 0$ for all $Y \in \hthree$ whenever $\rhoJ(X) = 0$. Fix
$X_0$ with $\rhoJ(X_0) = 0$: the polynomial
$Y \mapsto f(X_0, Y) = \sum_\alpha c_\alpha(X_0) \, m_\alpha(Y)$
is identically zero in $Y$, so linear independence of
$\{m_\alpha\}$ forces $c_\alpha(X_0) = 0$ for every $\alpha$. This
holds for every $X_0 \in V(\rhoJ)$, so each $c_\alpha$ vanishes on
$V(\rhoJ)$, hence lies in $I(V(\rhoJ)) = (\rhoJ)$ by (1a). Write
$c_\alpha(X) = \rhoJ(X) \cdot c_\alpha'(X)$ and resum:
\[
  f(X, Y) \;=\; \rhoJ(X) \cdot \sum_\alpha c_\alpha'(X) \, m_\alpha(Y)
  \;=:\; \rhoJ(X) \cdot g(X, Y),
\]
with $g \in \mathbb{R}[\hthree \times \hthree]$.

\emph{Inheriting invariance.} $\rhoJ$ is $\Ffour$-invariant on $X$
and constant in $Y$, so it is diagonal-$\Ffour$-invariant on the
pair. Combined with diagonal-$\Ffour$-invariance of $f$, the
identity $f = \rhoJ \cdot g$ forces $g$ to be diagonal-$\Ffour$-
invariant on $\{\rhoJ(X) \neq 0\} \times \hthree$, a Zariski-dense
open, and hence everywhere.

\emph{Degree bound.} $\rhoJ$ has matrix-entry degree $5$; if $f$
has minimal total matrix-entry degree $7$, $g$ has matrix-entry
degree $2$.

\textbf{Step 2 (degree-2 pair-invariants).} Set
\[
  a := \Tr X, \quad b := \Tr Y, \quad c := \Tr(X^2), \quad
  d := \Tr(X \cdot Y), \quad e := \Tr(Y^2).
\]
Under diagonal $\Ffour$, the degree-$2$ invariants on
$\hthree \times \hthree$ form the $6$-dimensional space
\[
  \mathcal{V}_2 \;=\; \mathrm{span}_{\mathbb{R}}\{\,a^2,\ ab,\ b^2,\ c,\ d,\ e\,\}.
\]
The first three are degree $2$ as products of degree-$1$ traces;
the last three are the degree-$2$ Faraut-Kor\'anyi pair-generators.
Diagonal $\Ffour$-invariance of $d = \Tr(X \cdot Y)$ is immediate
from $\Tr(gX \cdot gY) = \Tr(X \cdot Y)$ under the $\Ffour$-action
preserving the Jordan product.

\textbf{Step 3 (P1 kernel).} Setting $X = Y$ sends $(a, b, c, d, e)
\mapsto (a, a, c, c, c)$. The restriction map
$\mathcal{V}_2 \to \mathbb{R}[a, c]$ sends
$\{a^2, ab, b^2\} \mapsto a^2$ (each) and $\{c, d, e\} \mapsto c$
(each). The kernel of this restriction is $4$-dimensional, with
basis
\[
  k_0 = ab - a^2,\quad k_1 = b^2 - a^2,\quad k_2 = d - c,\quad k_3 = e - c,
\]
and a general element reads
\begin{equation}\label{eq:p1-kernel}
  g = \alpha\,(ab - a^2) + \beta'\,(b^2 - a^2) + \gamma\,(d - c) + \delta'\,(e - c),
\end{equation}
four real parameters $(\alpha, \beta', \gamma, \delta')$.

\textbf{Step 4 (non-negativity: block decomposition).} Every
$X \in \hthree$ decomposes uniquely as $X = \lambda_X I +
\mathring{d}_X + x_{\mathrm{off}}$, where
$\lambda_X = \tfrac{\Tr X}{3}$ (scalar block),
$\mathring{d}_X \in \mathbb{R}^3_0$ is the traceless diagonal
($\sum_i (\mathring{d}_X)_{ii} = 0$), and
$x_{\mathrm{off}} \in \mathbb{O}^3$ collects the three octonion
off-diagonal entries $x_{12}, x_{13}, x_{23}$. These three
summands are mutually orthogonal under the Jordan inner product
$\langle A, B \rangle_J = \Tr(A \cdot B)$: the scalar block
pairs with a traceless block to $\Tr(\lambda I \cdot
\mathring{d}) = \lambda \cdot \Tr(\mathring{d}) = 0$; the scalar
and off-diagonal pair to zero because the off-diagonal has zero
diagonal in $X$; and the traceless-diagonal and off-diagonal
pair to zero because $\mathring{d} \cdot x_{\mathrm{off}}$ has
zero diagonal. The factor of $2$ in
\[
  \Tr(X^2) \;=\; 3\lambda_X^2 + \lVert \mathring{d}_X \rVert^2 +
    2\lVert x_{\mathrm{off}} \rVert^2
\]
comes from the diagonal of $X^2$: entry
$(X^2)_{ii} = X_{ii}^2 + \sum_{j \neq i} |X_{ij}|_{\mathbb{O}}^2$
collects \emph{each} off-diagonal position $x_{ij}$ twice (once at
$i$, once at $j$) via
$X_{ij} \bar{X}_{ji} + X_{ji} \bar{X}_{ij} = 2|x_{ij}|^2$ (using
$X_{ji} = \bar{X}_{ij}$ for self-adjoint $X$). The bilinear form
$\Tr(X \cdot Y) = 3\lambda_X \lambda_Y +
\langle \mathring{d}_X, \mathring{d}_Y \rangle_{\mathbb{R}^3} +
2\,\mathrm{Re}\langle x_{\mathrm{off}}, y_{\mathrm{off}} \rangle_{\mathbb{O}^3}$
splits identically, with $\mathrm{Re}\langle a, b \rangle_{\mathbb{O}} =
\tfrac{1}{2}(a\bar{b} + b\bar{a})$ on each octonion slot.

Because diagonal $\Ffour$-invariance acts within each block
(scalar: trivially; traceless-diagonal: by the Weyl permutation
$S_3$; off-diagonal: by the induced octonionic rotation), each
block-level invariant of $(X, Y)$ is computed from its block
coordinates alone. Substituting the block expressions into
$g$ in~\eqref{eq:p1-kernel} and collecting by block, $g$ splits
additively as $g = g_A + g_B + g_C$, one summand per block, with
no cross-block interference. Each $g_*$ is a real quadratic form
in the pair of variables built from that block.

On Block~A (the scalar pair $(\lambda_X, \lambda_Y)$), the matrix
of the quadratic form in the basis $(\lambda_X^2,\ \lambda_X
\lambda_Y,\ \lambda_Y^2)$ — equivalently $(a^2, ab, b^2)$ up to
the factor $9$ from $a = 3\lambda_X$ — reduces to the $2 \times 2$
symmetric matrix
\[
  M_A \;=\; \begin{pmatrix} -(\alpha + \beta') & \alpha/2 \\ \alpha/2 & \beta' \end{pmatrix}
\]
in the coefficients of $(\lambda_X^2, \lambda_Y^2)$ with the
off-diagonal entry read off from the $\lambda_X \lambda_Y$
coefficient. Blocks B (traceless-diagonal pair
$(\mathring{d}_X, \mathring{d}_Y)$) and C (octonion off-diagonal
pair $(x_{\mathrm{off}}, y_{\mathrm{off}})$) each carry the same
structural quadratic form, parametrized by $(\gamma, \delta')$:
\[
  M_{BC} \;=\; \begin{pmatrix} -(\gamma + \delta') & \gamma/2 \\ \gamma/2 & \delta' \end{pmatrix}.
\]
Blocks B and C therefore share a PSD constraint.

\textbf{Step 5 (PSD constraints force $\alpha = -2\beta'$ and
$\gamma = -2\delta'$).} For $g \geq 0$ on all of
$\hthree \times \hthree$, each block matrix must be positive
semidefinite. PSD of $M_A$ requires $\beta' \geq 0$,
$-(\alpha + \beta') \geq 0$, and $\det M_A \geq 0$. A calculation
gives
\[
  4\,\det M_A \;=\; -\,(\alpha + 2\beta')^2,
\]
so $\det M_A \leq 0$ always, with equality iff
$\alpha = -2\beta'$. Combined with $\beta' \geq 0$, block~A forces
$\alpha = -2\beta'$, $\beta' \geq 0$, and Block~A reduces to
$\beta' \cdot (\Tr X - \Tr Y)^2$. Identically, $M_{BC}$ gives
$\gamma = -2\delta'$, $\delta' \geq 0$, with Blocks~B+C summing
to $\delta' \cdot \Jnorm{X - Y}^2$.

\textbf{Step 6 (P4 forces $\delta' > 0$).} Substituting
$\alpha = -2\beta'$, $\gamma = -2\delta'$ into~\eqref{eq:p1-kernel},
\[
  g(X, Y) \;=\; \beta'\,(\Tr X - \Tr Y)^2 \ +\ \delta'\,\Jnorm{X - Y}^2,
\]
a two-parameter non-negative cone with $\beta' \geq 0$ and
$\delta' \geq 0$. The $\beta'$-term depends only on $(\Tr X, \Tr Y)$
and is $\Tr(X \cdot Y)$-free; the $\delta'$-term, expanding
$\Jnorm{X - Y}^2 = \Tr(X^2) - 2\Tr(X \cdot Y) + \Tr(Y^2)$, has
$\Tr(X \cdot Y)$-coefficient $-2\delta'$. P4 (non-trivial
dependence on $\Tr(X \cdot Y)$) thus forces $\delta' > 0$.
Writing $\beta := \beta'$ and $\delta := \delta'$ gives the
statement. The SymPy script~\cite{phi-uniqueness} reproduces the
computation symbolically as an independent check.
\end{proof}

Theorem~\ref{thm:phi-family} leaves an Observable-reducible summand.
The $\beta$-term $(\Tr X - \Tr Y)^2$ is a polynomial in pointwise
invariants ($\Tr X$ and $\Tr Y$ are each in the single-copy invariant
ring), so it contributes a quantity that is \emph{fully recoverable}
from the Observable record of the Stream. The $\delta$-term, via the
pairwise invariant $\Tr(X \cdot Y)$ in its expansion
\eqref{eq:running-expansion}, is not.

This residual degree of freedom is resolved by adding path-space purity
as a sixth condition.

\begin{corollary}[Path-pure selection of $\Phiexp$]\label{cor:phi-canonical}
Under the additional requirement
\begin{itemize}[leftmargin=2em]
\item \textbf{(P6) Path-space purity.} $f$ carries no
Observable-reducible content - equivalently, $f$ has no summand
expressible as a polynomial in the pointwise $\Ffour$-invariants
$(\Tr X, \Tr X^2, \det X)$ and
$(\Tr Y, \Tr Y^2, \det Y)$ alone,
\end{itemize}
the family \eqref{eq:phi-family} collapses to $\beta = 0$, giving
\begin{equation}\label{eq:phi-canonical}
  \boxed{\;\Phiexp[\Stream] \;=\; \delta \sum_k \rhoJ(X_k)\,\Jnorm{X_{k+1} - X_k}^2\;}
\end{equation}
unique up to a positive scalar $\delta > 0$.
\end{corollary}

(P6) is not a further selection condition of the same type as
(P1)-(P5); it is a projection onto the non-Observable-reducible
summand of the minimal-degree family. Writing $\Phiexp$ without this
projection would include a $\beta \geq 0$ coefficient on a
polynomial in pointwise $\Ffour$-invariants that an outside party
can recover by measuring Kinds alone - content already visible on
the Observable side. The path-pure $\Phiexp$ above is the reduced
representative of the family; we fix the remaining scalar freedom
by setting $\delta = 1$.

% ====================================================================
% SECTION 5: STRONG-FORM COMPLEMENTARITY
% ====================================================================

\section{Strong-form complementarity: the Observable-Beable split}
\label{sec:complementarity}

Observable extraction and $\phi$-iteration on an $\hthree$-Stream do
not commute. The claim splits by regime. In the autonomous case
($\phi$ $\Ffour$-equivariant, gauge fixed by the dynamics) the
non-commutation reduces to a single residual global gauge
(Theorem~\ref{thm:autonomous-comp}); the autonomous setting is thus
the weakest, and autonomous $\Phiexp$ is Observable-reducible. In
the non-autonomous case - a self-modeler whose $\phi$-step is
modulated by an external drive, so that the $\Ffour$-frame the
dynamics uses at step $k$ is not predetermined by the state
$X_k$ alone - the pairwise invariant $\Tr(X_k \cdot X_{k+1})$
escapes the pointwise Kind sequence; strong-form complementarity
follows unconditionally (Proposition~\ref{prop:strong-comp}). The
\emph{non-autonomous} here is the analogue of a driven Hamiltonian
in QM: the algebra supplies the state space and the step rule, an
external source chooses the gauge-varying correction; no coupling
to a named heat bath is assumed, only that $\phi$ fails to be
$\Ffour$-equivariant. The resulting Observable-Beable exhaustion
(Corollary~\ref{cor:OB1}) says the two information modes jointly
span the diagonal $\Ffour$-invariant content of the Stream; it
realizes, inside $\hthree$ by algebra, Bell's observable/beable
split. We use ``complementarity'' in Bell's
observable/beable sense (two modes exhaust the available
information), not in Bohr's measurement-context sense (two
mutually exclusive preparations).

\subsection{The mechanism}
\label{sec:complementarity-mechanism}

Paper~5 selects, under van~de~Wetering's sequential-product
axioms~\cite{vdw2019b,alfsen-shultz2001}, the Alfsen-Shultz
sequential product $X * Y = \sqrt{X}\, Y\, \sqrt{X}$ as the update
taking state-before-measurement to state-after-measurement on any
Euclidean Jordan algebra the framework admits; $X * Y \neq Y * X$ in
general. The non-commutation that drives the results below is a
different one - between Observable projection and $\phi$-iteration,
not between sequential products of a pair of measurements - but
both live inside the same algebra.

Two operations are available to a Stream segment $(X_k, X_{k+1})$.

\begin{itemize}[leftmargin=2em]
\item The \textbf{Observable projection}
$P_\mathcal{O} : X \mapsto [X] \in \Omega / \Ffour$, which maps a
27-dimensional State to its 3-dimensional Kind; equivalently, reads off
the three invariants $(\Tr X, \Tr(X^2), \det X)$.
\item The \textbf{Stream propagator} $\phi : X_k \mapsto X_{k+1}$, the
self-modeling iteration of Section~\ref{sec:stream}.
\end{itemize}

These do not commute in a precise and testable sense. Applying $\phi$
first and then $P_\mathcal{O}$ gives a well-defined composition:
$P_\mathcal{O}(\phi(X_k)) = [X_{k+1}]$. Applying $P_\mathcal{O}$ first
returns the orbit $[X_k]$; the next $\phi$ application requires a
\emph{choice of gauge representative}, because $\phi$ operates on States
and $[X_k]$ is not a State. Any such choice injects gauge information
that was destroyed in the projection, and different choices produce
different propagations. The commutator
$[\phi, P_\mathcal{O}]$ is not a bounded operator; it is a statement
about well-definedness. Observable extraction destroys the 24
gauge-dimensional degrees of freedom that the next $\phi$ step would have
consumed.\footnote{Generic $\Ffour$-orbits on $\hthree$ have dimension
$24$: the spectrum of an element supplies $3$ independent
$\Ffour$-invariants ($\Tr$, $\Tr(X^2)$, $\det$) on the $27$-dimensional
ambient, so the complementary orbit direction is $27 - 3 = 24$. The
dimension count agrees with $\dim \Ffour - \dim \mathrm{Spin}(8) =
52 - 28 = 24$, with $\mathrm{Spin}(8)$ the generic stabilizer. See
\cite{harvey1990} \S14.}

\subsection{Autonomous case: the weakest setting}
\label{sec:autonomous-complementarity}

\begin{theorem}[Autonomous complementarity]\label{thm:autonomous-comp}
Let $\Stream$ be a Stream of a self-modeler in $\hthree$ whose
self-modeling operator $\phi$ is $\Ffour$-equivariant and whose dynamics
are autonomous (gauge fixed throughout by $\phi$ itself). Then:
\begin{itemize}[leftmargin=2em]
\item \textbf{(i)} The pointwise Kind sequence $([X_0], [X_1], \ldots)$
determines the Stream up to a single global $\Ffour$-gauge.
\item \textbf{(ii)} Any polynomial Observable-side functional of the
Kind sequence carries strictly less information than the Stream itself:
the gauge orbit of $\Stream$ is the maximal Observable-reducible content.
\item \textbf{(iii)} $\Phiexp[\Stream]$ is Observable-reducible in
the autonomous case: its non-Observable content reduces to a single
global $\Ffour$-gauge orbit, which is unphysical.
\end{itemize}
\end{theorem}

\begin{proof}[Proof sketch]
For \textbf{(i)}, fix $X_0 \in \Omega$ as a representative of its
Kind $[X_0]$. By Faraut-Kor\'anyi, every $\Ffour$-invariant of $X_0$
factors through $[X_0]$, so any two representatives of $[X_0]$
differ by an element of $\Ffour$. $\Ffour$-equivariance of $\phi$
gives $\phi(g X_0) = g \phi(X_0)$; iterating, $\phi^k(g X_0) =
g \phi^k(X_0)$ for all $k$. The whole Stream thus propagates with a
single global gauge $g$ applied to the whole trajectory, and the
Kind sequence determines the Stream up to that single global choice.

For \textbf{(ii)}, a polynomial functional of the Kind sequence is a
polynomial in $\{\Tr X_k, \Tr(X_k^2), \det X_k\}_k$. By (i) and
Faraut-Kor\'anyi, this ring coincides with the diagonal-$\Ffour$-invariant
ring on the Stream modulo the single global gauge. The Stream itself
sits in $\Omega^{\mathbb{N}}$, one global-gauge orbit above the Kind
sequence; the gauge orbit is the maximal Observable-recoverable
content.

For \textbf{(iii)}, $\Phiexp[\Stream]$ depends only on the data
$(\rhoJ(X_k), \Jnorm{X_{k+1} - X_k}^2)$, both of which are
$\Ffour$-invariant under the diagonal action. So
$\Phiexp[g \cdot \Stream] = \Phiexp[\Stream]$ for any
$g \in \Ffour$. Under (i), the Kind sequence plus a global gauge
determines the Stream; since $\Phiexp$ does not see the global gauge,
$\Phiexp$ is a functional of the Kind sequence alone and thus
Observable-reducible.
\end{proof}

\subsection{Non-autonomous complementarity}
\label{sec:non-autonomous-complementarity}

In the non-autonomous case - drive variations across phi-time - the
situation is sharper.

Consider two Streams $\Stream_a$ and $\Stream_b$ constructed so that their
pointwise Kind sequences match exactly,
$[X_k^{(a)}] = [X_k^{(b)}]$ for all $k$, but whose gauge representatives
are transformed by a time-varying element $g_k \in \Ffour$ between them.
In the autonomous setting the two Streams differ by a global gauge and
carry identical $\Ffour$-invariant content. In the non-autonomous setting
- when $g_k$ is not constant across $k$ - the pairwise invariant
$\Tr(X_k \cdot X_{k+1})$ can differ between the two Streams, because
\[
  \Tr(g_k X_k^{(a)} \cdot g_{k+1} X_{k+1}^{(a)})
  \;\neq\; \Tr(X_k^{(a)} \cdot X_{k+1}^{(a)})
  \quad \text{in general},
\]
whereas the pointwise Kinds are preserved. The pair-level invariant
escapes the pointwise Kind sequence.

The explicit $21\times$ demonstration of
Section~\ref{sec:phi-properties} realizes this: the ratio is
$\Phiexp$-content of the path-space residual, invisible to any
outside party measuring only pointwise Kinds.

\begin{proposition}[Strong-form complementarity, non-autonomous]\label{prop:strong-comp}
For non-autonomous Streams in $\hthree$ with time-varying $\Ffour$-gauge,
the pointwise Kind sequence does \emph{not} determine the diagonal
$\Ffour$-invariant ring on path-space. In particular, the pairwise
invariant $\Tr(X_k \cdot X_{k+1})$ belongs to a strict super-ring of
the pointwise invariant ring on $\hthree^N$, and polynomial diagonal
$\Ffour$-invariant functionals of the Stream that depend on
path-space invariants are strictly richer than those that depend on
pointwise Kinds alone.
\end{proposition}

\begin{proof}[Proof]
At $N = 2$, take $X_0 = X_0' = \mathrm{diag}(1/2, 3/10, 1/5)$ and
$X_1 = \mathrm{diag}(2/5, 7/20, 1/4)$,
$X_1' = \sigma \cdot X_1 = \mathrm{diag}(7/20, 2/5, 1/4)$, where
$\sigma$ is the transposition of the first two diagonal entries.
$\sigma$ is a Jordan-algebra automorphism of $\hthree$ (permuting
the three primitive idempotents preserves the Jordan product), and
thus belongs to $\Ffour = \Aut\hthree$. All four states lie in the
interior $\Omega^\circ$ and share pointwise $\Ffour$-invariants:
$\Tr X_0 = \Tr X_1 = \Tr X_1' = 1$,
$\Tr(X_0^2) = 19/50$, $\Tr(X_1^2) = \Tr(X_1'^2) = 69/200$,
$\det X_0 = 3/100$, $\det X_1 = \det X_1' = 7/200$, so any
polynomial in pointwise $\Ffour$-invariants takes equal values on
the two pairs $(X_0, X_1)$ and $(X_0', X_1')$. The weights are
strictly positive: $\rhoJ(X_0) = 7/5000$,
$\rhoJ(X_1) = \rhoJ(X_1') = 49/120000$. But
$\Tr(X_0 \cdot X_1) = 71/200$ while $\Tr(X_0' \cdot X_1') = 69/200$,
and $\Tr(X \cdot Y)$ is invariant under the \emph{diagonal}
$\Ffour$-action $(X, Y) \mapsto (gX, gY)$ with the \emph{same}
$g$ applied to both entries. Hence $\Tr(X \cdot Y)$ is not in the
pointwise ring; the super-ring is strict at $N = 2$. Since
$\rhoJ > 0$ on both streams, the witness doubles as a $\Phiexp$
discriminator:
$\Phiexp[(X_0, X_1)] = \rhoJ(X_0) \cdot \Jnorm{X_0 - X_1}^2 =
21/10^6$ while $\Phiexp[(X_0', X_1')] = 49/10^6$, a $7{:}3$ ratio
between Streams with identical pointwise Kinds and all Observable
weights positive. For any $N > 2$, apply the witness at any
consecutive pair $(k, k+1)$ with the remaining indices held fixed:
the pointwise Kind sequence matches while the pair-invariant at
$(k, k+1)$ does not. The richness claim follows: $\Phiexp$ itself,
by (P3) and (P4) of Section~\ref{sec:phi-properties}, is a
polynomial diagonal $\Ffour$-invariant functional that
distinguishes Streams with matched pointwise Kinds and nonzero
$\rhoJ$-weight.
\end{proof}

\begin{remark}[Depth of the path-space residual]\label{rem:path-space-depth}
Proposition~\ref{prop:strong-comp} witnesses the super-ring at the
pair level. The $\Phiexp$ of
Definition~\ref{def:phi} is built as a sum of pair-local terms
$\rhoJ(X_k) \Jnorm{X_{k+1} - X_k}^2$, so only pair-level diagonal
$\Ffour$-invariants enter its definition; triple and higher
diagonal invariants are out of scope for the uniqueness result of
Corollary~\ref{cor:phi-canonical}, which is explicitly a pair-level
statement. An explicit finite generating set for the full diagonal
$\Ffour$-invariant ring on $\hthree^N$ is open
(Section~\ref{sec:gaps}, G2); Faraut-Kor\'anyi Chapter~V gives the
invariant ring for a single copy, but the diagonal product case is
not in the classical literature.
\end{remark}

\subsection{Observable-Beable exhaustion}
\label{sec:OB1}

\begin{corollary}[Observable-Beable exhaustion]\label{cor:OB1}
Two modes of information about a Stream are available within the
four-type ontology. The \emph{Observable mode} ($\mathcal{O}$):
any $\Ffour$-invariant polynomial functional of the pointwise Kind
sequence; by Proposition~\ref{prop:strong-comp} this mode is
strictly weaker than the full diagonal $\Ffour$-invariant ring on
path-space. The \emph{Beable mode} ($\mathcal{B}$): direct
specification of the Stream as an element of
$\Omega^{\mathbb{N}} / \Ffour$, recovering the pair-level invariants
$\Tr(X_k \cdot X_{k+1})$ that $\mathcal{O}$ misses. Observable
extraction consumes the $\Ffour$-gauge information that the next
$\phi$-iteration would use, so no $\Ffour$-invariant refinement of
$\mathcal{O}$ closes the gap in the non-autonomous regime; the two
modes jointly span the diagonal $\Ffour$-invariant content of the
Stream. This is Bell's~\cite{bell1984speakable} observable/beable
split realized inside $\hthree$ by algebra.
\end{corollary}

% ====================================================================
% SECTION 6: PROPERTIES OF PHI
% ====================================================================

\section{Properties of $\Phiexp$}
\label{sec:consequences}

Temporal additivity is a direct consequence of
Definition~\ref{def:phi}'s sum form. Spatial aggregation sits
outside the uniqueness result and is flagged as an open question
(Section~\ref{sec:gaps}, G3); we do not state a conjectured form
here.

\subsection{Temporal compositionality (theorem)}
\label{sec:temporal-composition}

\begin{proposition}[Temporal additivity]\label{prop:temporal-add}
For a Stream $\Stream = (X_0, \ldots, X_K)$ and any index $0 \leq k^* \leq K$,
\begin{equation}\label{eq:temporal-add}
  \Phiexp[(X_0, \ldots, X_K)]
  \;=\; \Phiexp[(X_0, \ldots, X_{k^*})]
  \;+\; \Phiexp[(X_{k^*}, \ldots, X_K)].
\end{equation}
\end{proposition}

\noindent
Immediate from the sum form.

% ====================================================================
% SECTION 7: EMPIRICAL SIGNATURES
% ====================================================================

\section{Empirical signatures}
\label{sec:signatures}

A measure that lives inside $\hthree$ and is by construction
Observable-indistinguishable for pairs of Streams with matching
pointwise Kinds can still be tested from outside. Three classes of
signature do the work.

\paragraph{(S1) Matched-pair discriminator at $21\times$.}
The construction of Section~\ref{sec:phi-properties} provides a
clean class-level discriminator: two Streams with pointwise Kinds
matched to machine precision and $\Phiexp$ differing by a factor
of $21$. Any \emph{theory} claiming that a Stream's $\Phiexp$ (or
an analogous Beable-side functional) is recoverable from pointwise
$\Ffour$-invariants alone must report the pair as $\Phiexp$-equal;
the pair falsifies that theory class. Any apparatus accessing the
pair-level invariant $\Tr(X_k \cdot X_{k+1})$ distinguishes them.
The test runs unconditionally in $\hthree$-arithmetic; it does not
require a biological or physical encoding. Application to a
specific target system is a separate question of whether that
system \emph{admits} matched-pair preparation at all.

\paragraph{(S2) Depth-scaling (conjectural).}
The per-tick contribution $\rhoJ(X_k)\Jnorm{X_{k+1} - X_k}^2$ is a
degree-$7$ polynomial in the entries of $X_k$ and $X_{k+1}$;
under a family of nested self-models indexed by a depth parameter,
iterating at increasing depth should drive a polynomial-bounded
amplification of $\Phiexp$ whose leading order is set by this
degree structure. Converting this algebraic expectation into a
measured profile requires an encoding (G1); without one, (S2) is
conjectural motivation, not a measurement protocol. The prediction,
conditional on any fixed encoding: \emph{some} polynomial-bounded
amplification profile obtains with depth, and flat or exponential
scaling would falsify the depth-scaling form.

\paragraph{(S3) Algebraic zero class.}
$\Phiexp = 0$ exactly on two algebraic classes of Stream: (i) static
Streams ($X_{k+1} = X_k$, by P1), and (ii) Streams confined to the
$\rhoJ$-vanishing locus (by P2). Within a fixed encoding, systems
whose dynamical image lies in these classes - static configurations
or configurations on the rank-deficient / maximally-mixed stratum -
carry $\Phiexp = 0$ at the algebraic level. The identification of
real-world exemplars with these classes (which real systems, if
any, encode into the $\rhoJ = 0$ locus under a biological encoding)
is an encoding question (G1), not an algebraic one.

% ====================================================================
% SECTION 8: RELATED WORK
% ====================================================================

\section{Related work}
\label{sec:related}

Three neighbours bear mention: IIT (as the other $\Phi$ with a
consciousness-adjacent target), Barandes's indivisible-stochastic
reformulation of QM (as the source of the trajectory-as-beable
shape), and the Jordan-algebraic reconstruction programs (as the
upstream machinery Paper~5 uses). Each is different from this
paper in a specific way, and these differences - not the surface
resemblance - are where the content is.

\paragraph{IIT.} Integrated Information Theory
\cite{oizumi2014,tononi2016,albantakis2023} shares an initial letter
with $\Phiexp$ and a motivating target. The two differ in substrate
(partitioned discrete physical system versus $\hthree$-valued iterates),
in derivational grounding (IIT's five axioms versus the uniqueness
theorem on Streams of Section~\ref{sec:uniqueness}), and in compositional
predictions (IIT's integration and exclusion postulates make spatial
decomposition destructive; $\Phiexp$ is temporally additive by
Proposition~\ref{prop:temporal-add}, with spatial behaviour open).
They are different objects with similar names. A fuller
discriminator table is below.

\begin{table}[H]
\centering
\small
\begin{tabular}{@{}p{0.25\textwidth}p{0.32\textwidth}p{0.32\textwidth}@{}}
\toprule
\textbf{Dimension} & \textbf{IIT 4.0} & \textbf{This paper's $\Phiexp$} \\
\midrule
Substrate & Partitioned discrete physical system; cause-effect structure over state transitions & $\hthree$-valued iterates; algebraic self-modeling on a Euclidean Jordan algebra \\
Selection & Five phenomenal axioms (existence, intrinsicality, information, integration, exclusion) + their physical postulates & Observable-Beable ontology (Def~\ref{def:ontology}); six conditions on pair-invariants (P1-P6), minimal matrix-entry degree, non-negativity \\
Uniqueness & Measure form derived from the five phenomenal axioms applied to cause-effect structures; the minimum-information-partition argument then selects \emph{which} subsystem is the conscious complex & Finite-dimensional enumeration of diagonal $\Ffour$-invariants on pairs (Thm~\ref{thm:phi-family}) + path-space purity projection (Cor.~\ref{cor:phi-canonical}) \\
Temporal composition & Transition-local evaluation; no closed additivity statement & Temporally additive (Prop.~\ref{prop:temporal-add}) \\
Spatial composition & Exclusion postulate: one maximally integrated complex at each scale, others suppressed & Open (Section~\ref{sec:gaps}, G3); no conjectured form stated in this paper \\
\bottomrule
\end{tabular}
\end{table}

A direct empirical discriminator requires a matched protocol -
varying metacognitive depth or a Stream-level gauge schedule - not
yet implemented; Section~\ref{sec:signatures} (S2) records the
depth-scaling prediction as the cleanest candidate.

\paragraph{Barandes's indivisible-stochastic framework.}
Barandes~\cite{barandes2023stochquantum} reformulates quantum
mechanics as an indivisible stochastic process on an underlying
sample-trajectory ontology, with the Hilbert-space wave function
serving as bookkeeping. The trajectory-beable / Observable-record
split in Section~\ref{sec:ontology} adopts this shape; the
present framework uses $\hthree$-valued iterates where Barandes
uses discrete sample paths, and does not import the stochastic
machinery.

\paragraph{Reconstruction programs for QM.}
Van de Wetering's sequential-product axiomatization of Euclidean
Jordan algebras~\cite{vdw2019a,vdw2019b}, Alfsen-Shultz state-space
geometry~\cite{alfsen-shultz2001}, and the homogeneous-self-dual-cone
approach to QM reconstruction give the upstream mathematical setting
Paper~5~\cite{ehrlich2026qm} uses to obtain $\hthree$. The measure
$\Phiexp$ is one step past the reconstruction: having obtained the
algebra and its dynamics, it asks which $\Ffour$-invariant functional
on trajectories is selected by non-reducibility and path-space
purity.

% ====================================================================
% SECTION 9: OPEN GAPS
% ====================================================================

\section{Open gaps}
\label{sec:gaps}

We register the principal open items with their status and minimal
next test.

\paragraph{(G1) Biological encoding (framework task, not theorem gap).}
The theorem and the measure are encoding-neutral at the $\Ffour$-orbit
level. Which particular $\hthree$-encoding corresponds to the state
variables of a biological nervous system is an independent empirical
question. The paper's central claims do not depend on fixing an
encoding; applying $\Phiexp$ numerically to a specific neuronal
recording does. Cheapest test: pick a candidate encoding motivated by a
Markov-blanket or free-energy-principle model and compute
$\Phiexp$ on an existing recording.

\paragraph{(G2) Explicit generators for the diagonal
$\Ffour$-invariant ring on $\hthree^N$.}
Proposition~\ref{prop:strong-comp} shows the diagonal
$\Ffour$-invariant ring on $\hthree^N$ strictly contains the
pointwise invariant ring, but an explicit finite generating set is
not in the classical literature. Faraut-Kor\'anyi Chapter V gives
the $N=1$ ring. Cheapest advance: compute a generating set at
$N=2$.

\paragraph{(G3) Spatial compositionality.}
Temporal additivity is a theorem (Prop.~\ref{prop:temporal-add});
spatial aggregation is open. The question has two sub-gaps: (a)
what counts as a ``subsystem'' of a self-modeler in $\hthree$,
given that $\hthree$ is simple and does not decompose as a Jordan
direct sum; (b) how the $\Phiexp$ of a composite self-modeler -
presumably living in a larger Jordan or C*-algebra - relates to
the $\Phiexp$'s of its parts. Neither sub-gap has a natural
conjectured form at this stage. Cheapest advance: fix a composite
algebra (e.g., $\hthree \otimes \hthree$ as a Jordan algebra via
the van~de~Wetering composite construction) and compute $\Phiexp$
on a two-subsystem Stream with a shared input, comparing against
$\Phiexp_1 + \Phiexp_2$.

\paragraph{(G4) Unconditional $\rhoJ$ uniqueness
(Remark~\ref{rem:rho-alternatives}).} Theorem~\ref{thm:rho-uniqueness}
assumes the factored shape $p \cdot q$ at unbounded polynomial degree;
the Lean theorem \texttt{rhoJ\_unique\_minimal}~\cite{lean-rhoj} drops
factorization at eigenvalue degree $\leq 5$. The unconditional
statement - every $\Ffour$-invariant polynomial on $\hthree$ vanishing
on $\partial\Omega$ and at $[I/3]$ with $\rhoJ$-minimal degree is
$c \cdot \rhoJ$ - remains open. Cheapest advance: extend the Lean
proof to unbounded degree, or give a combined factorization-free
argument at the ambient eigenvalue ring.

\paragraph{(G5) Self-measurement inaccessibility (research direction).}
Corollary~\ref{cor:OB1} (Observable-Beable exhaustion) is the \emph{weak}
form: an outside party reading pointwise Kinds misses the path-space content
$\Phiexp$ carries. The \emph{strong} form - that no self-modeler can recover
$\Phiexp$ of its \emph{own} Stream from any Observable record it can persist -
is a research target: the algebraic analogue of a G\"odel/Tarski/Lawvere
self-reference limit, with a no-cloning argument supplying the
``cannot persist a copy of the prior frame'' step. Stated, unproven; see the
program note \texttt{research/phi-inaccessibility-program.md}. If closed, it
would place the inaccessibility of $\Phiexp$ on the footing of Tarski's
undefinability of truth: real in the algebra, unrecoverable from inside.

\bigskip

\noindent
The paper's status is asymmetric. Autonomous and non-autonomous
complementarity, the $\rhoJ$ selection (modulo G4), the $\Phiexp$
pair-family at minimal matrix-entry degree, the path-pure
projection, Observable-Beable exhaustion, and temporal additivity
are theorems; what remains open is explicit generators for the
diagonal $\Ffour$-invariant ring on $\hthree^N$ (G2), spatial
aggregation (G3), the factorization-free $\rhoJ$ uniqueness (G4),
and the biological encoding question deliberately not committed
to (G1). The autonomous case is a theorem; the non-autonomous
case is the falsifier program; the uniqueness half stands
unconditionally modulo G4. The matched-pair $21\times$
discriminator sits inside the algebra and requires none of G1-G4
to run.

% ====================================================================
% BIBLIOGRAPHY
% ====================================================================

\begin{thebibliography}{99}

\bibitem{alfsen-shultz2001}
E.~M.~Alfsen and F.~W.~Shultz.
\newblock \emph{State Spaces of Operator Algebras: Basic Theory, Orientations, and $C^*$-Products}.
\newblock Birkh\"auser, 2001.

\bibitem{barnum-graydon-wilce2015}
H.~Barnum, M.~A.~Graydon, and A.~Wilce.
\newblock Composites and categories of {E}uclidean {J}ordan algebras.
\newblock arXiv:1606.09331, 2015.

\bibitem{bell1984speakable}
J.~S.~Bell.
\newblock \emph{Speakable and Unspeakable in Quantum Mechanics}.
\newblock Cambridge University Press, 2nd ed., 2004 (orig.\ 1987).

\bibitem{barandes2023stochquantum}
J.~A.~Barandes.
\newblock The stochastic-quantum correspondence.
\newblock \emph{Philosophy of Physics}, 3(1):8, 2025.
\newblock arXiv:2302.10778.

\bibitem{bengtsson-zyczkowski2006}
I.~Bengtsson and K.~\.Zyczkowski.
\newblock \emph{Geometry of Quantum States}.
\newblock Cambridge University Press, 2006.

\bibitem{ehrlich2026qm}
B.~Ehrlich.
\newblock Complex quantum mechanics from faithful self-modeling (Paper~5).
\newblock Submitted to \emph{Journal of Mathematical Physics}, 2026.

\bibitem{ehrlich2026h3o}
B.~Ehrlich.
\newblock The exceptional Jordan algebra as the unique non-composable ground for self-modeling (Paper~7).
\newblock \emph{ehrlich.dev/papers/sm-from-self-modeling/main.pdf}, 2026.

\bibitem{devaney1989}
R.~L.~Devaney.
\newblock \emph{An Introduction to Chaotic Dynamical Systems}.
\newblock Addison-Wesley, 2nd ed., 1989.

\bibitem{faraut-koranyi1994}
J.~Faraut and A.~Kor\'anyi.
\newblock \emph{Analysis on Symmetric Cones}.
\newblock Oxford University Press, 1994.

\bibitem{harvey1990}
F.~R.~Harvey.
\newblock \emph{Spinors and Calibrations}.
\newblock Academic Press, 1990.

\bibitem{hofstadter1979}
D.~R.~Hofstadter.
\newblock \emph{G\"odel, Escher, Bach: An Eternal Golden Braid}.
\newblock Basic Books, 1979.

\bibitem{hofstadter2007}
D.~R.~Hofstadter.
\newblock \emph{I Am a Strange Loop}.
\newblock Basic Books, 2007.

\bibitem{katok-hasselblatt1995}
A.~Katok and B.~Hasselblatt.
\newblock \emph{Introduction to the Modern Theory of Dynamical Systems}.
\newblock Cambridge University Press, 1995.

\bibitem{phi-numerics}
B.~Ehrlich.
\newblock Numerical verification of $\Phi$ properties P1-P5 on $h_3(\mathbb{O})$.
\newblock \emph{research/qualia-fixed-point/phenomenological\_measure.py}, 2026.

\bibitem{lean-rhoj}
B.~Ehrlich.
\newblock Machine-verified uniqueness of $\rho_J$ at eigenvalue degree $\leq 5$.
\newblock \emph{Lean source: RadicalRelativity/RhoJ.lean, theorem \texttt{rhoJ\_unique\_minimal}}, 2026.

\bibitem{phi-uniqueness}
B.~Ehrlich.
\newblock SymPy-verified enumeration of the minimal-degree $\Phi$-family on $h_3(\mathbb{O})$.
\newblock \emph{research/qualia-fixed-point/phi\_uniqueness.py}, 2026.

\bibitem{springer-veldkamp2000}
T.~A.~Springer and F.~D.~Veldkamp.
\newblock \emph{Octonions, Jordan Algebras and Exceptional Groups}.
\newblock Springer, 2000.

\bibitem{tononi2016}
G.~Tononi et~al.
\newblock Integrated information theory: from consciousness to its physical substrate.
\newblock \emph{Nature Reviews Neuroscience}, 2016.

\bibitem{oizumi2014}
M.~Oizumi, L.~Albantakis, and G.~Tononi.
\newblock From the phenomenology to the mechanisms of consciousness: Integrated Information Theory 3.0.
\newblock \emph{PLOS Computational Biology}, 10(5):e1003588, 2014.

\bibitem{albantakis2023}
L.~Albantakis, W.~Barbosa, G.~Findlay, M.~Grasso, A.~M. Haun, W.~Marshall, W.~G.~P. Mayner, A.~Zaeemzadeh, M.~E. Boly, B.~E. Juel, S.~Sasai, K.~Fujii, I.~David, J.~Hendren, J.~P. Lang, and G.~Tononi.
\newblock Integrated Information Theory (IIT) 4.0: Formulating the properties of phenomenal existence in physical terms.
\newblock \emph{PLOS Computational Biology}, 19(10):e1011465, 2023.

\bibitem{vdw2019a}
J.~van de Wetering.
\newblock An effect-theoretic reconstruction of quantum theory.
\newblock \emph{Compositionality}, 1:1, 2019.

\bibitem{vdw2019b}
J.~van de Wetering.
\newblock Sequential product spaces are {J}ordan algebras.
\newblock \emph{J.~Math.\ Phys.}, 60(6):062201, 2019.

\bibitem{weyl1939}
H.~Weyl.
\newblock \emph{The Classical Groups: Their Invariants and Representations}.
\newblock Princeton University Press, 1939.

\end{thebibliography}

% ====================================================================
% APPENDIX A: STANDALONE SELF-MODELING DEFINITION
% ====================================================================

\appendix

\section{Self-modeling system (standalone)}
\label{app:self-modeling}

This appendix reproduces Paper~5~\cite{ehrlich2026qm}'s definition
of a self-modeling system for readers who prefer not to follow the
pointer. The definition is four-clause; the theorem of Paper~5
(Thm.~6.3 there) shows that every finite-dimensional self-modeling
system has state space $\mathbb{R}^n$- or $M_n(\mathbb{C})^{\mathrm{sa}}$-like
structure, with the non-commutative case selected by local
tomography. For the present paper, only the definition is used; it
fixes $(\hthree, \Ffour, \Omega, \phi)$ as the algebra, group,
state space, and normalized-square iteration of interest.

\begin{definition}[Self-modeling system, Paper~5 Def.~2.6]
\label{app-def:sms}
A \emph{self-modeling system} is a tuple $(V, \varphi, V_{BM})$
satisfying:
\begin{enumerate}[label=(\roman*)]
  \item $V$ is a nontrivial finite-dimensional spectral order-unit
    space: $V$ has at least two orthogonal nontrivial projective
    units.
  \item $\varphi: V_B \to V_M$ is an order isomorphism between the
    body component $V_B$ and the model component $V_M$
    (\emph{faithful tracking}).
  \item $V_{BM}$ is the minimal composite order-unit space carrying
    product states, product effects, non-signaling constraints, and
    a product-form sequential product (\emph{minimal internal
    composite}).
  \item $V$ has no nontrivial direct-sum decomposition
    $V = V_1 \oplus V_2$ into nonzero order-unit subspaces
    (\emph{simplicity}).
\end{enumerate}
\end{definition}

The four clauses unpack the three words of the concept: (i) = a
system, (ii) = modeling (faithful internal tracking), (iii) = self-
(internal, minimal composite), (iv) = the whole (no hidden
sectors). Paper~5 shows this implies, via van~de~Wetering's
sequential-product theorem and Artin-Wedderburn, that $V$ is
$M_n(\mathbb{C})^{\mathrm{sa}}$ in the complex case, with the
normalized iteration $\phi(X) = X^2 / \Tr(X^2)$ emerging as the
sequential-product self-application at the maximally coherent
face. Paper~7~\cite{ehrlich2026h3o} then selects $\hthree$ as the
unique non-composable instance, which is the setting of the
present paper.

\end{document}
